% !TeX root = ../教材题库_全部条目_转写工作稿.tex
% Source: XJ-B1-C05-EX-P157-Y0962
% 第5章 / 5.1 任意角和弧度制 / 教材 p.157 / PDF 第164页
% OCR未校对：公式与文字均可能有误；未完成数学TeX转写；禁止直接发给学生
% 原题及相邻原解截图（核对用）：../题目原图/XJ-B1-C05-EX-P157-Y0962-01.webp
% 本文件只在转写完成后作为题干使用；OCR 条目尚不是可靠题干。
\begin{bankitem}{5.1 例题 例题1}{第5章 / 5.1 任意角和弧度制 / 教材 p.157 / PDF 第164页}{XJ-B1-C05-EX-P157-Y0962}{OCR未校对：公式与文字均可能有误；未完成数学TeX转写；禁止直接发给学生}
\begin{ocrdraft}
[教材p.157 / PDF第164页]
\par
00在0。、360。内找出与下列各角终边相同的角，并判定它是第几象限角．
\par
（2）1600。．
\par
（3）
\par
一819。36'．
\par
分析只需将这些角表示成仅十差\ensuremath{\cdot}360。（0。a<360。）的形式，然后根据角来
\par
确定它们所在的象限．
\par
解（1)因为一80。一280。一360。，
\par
所以在0。、360。内，与一80。角终边相同的角是280。，它是第四象限角．
\par
（2）因为1600。一160。十4\ensuremath{\times}360。
\par
所以在0。、360。内，与1600。角终边相同的角是160。，它是第二象限角．
\par
（3）因为一819。36'一260。24'一3\ensuremath{\times}360。
\par
所以在0。、360。内，与一819。36/角终边相同的角是260。24'，它是第三象限角．
\end{ocrdraft}
\end{bankitem}
